% !TeX root = ../main.tex
\chapter{知识库}
\label{sec:knowledge-guardrail}

本章介绍 OralSage 从多源医学资料到智能体知识服务的完整构建链路。教材、临床规范、诊疗手册和题库经过结构恢复与语义处理后，分别形成医学文本、知识图谱和题目对象；这些对象通过多路索引与来源关系连接，再由 KAGSolver 为课程规划、资源生成、教学讲解、智能答疑和学习测评提供带有知识关系与来源位置的检索上下文。

\section{数据输入与预处理}

本项目收集了14本口腔医学专业课程教材、口腔医学领域专家共识与临床规范、《默沙东诊疗手册》以及口腔医师执业考试真题作为知识库基础资料输入，具体资料说明如表~\ref{tab:kb-source-overview}所示。

{\small
\renewcommand{\arraystretch}{1.08}
\begin{longtable}{p{0.27\linewidth}p{0.63\linewidth}}
  \caption{知识库基础资料}\label{tab:kb-source-overview}\\
  \toprule
  资料类型 & 内容与使用范围 \\
  \midrule
  \endfirsthead
  \multicolumn{2}{c}{续表~\thetable}\\
  \toprule
  资料类型 & 内容与使用范围 \\
  \midrule
  \endhead
  \bottomrule
  \endfoot
  口腔医学教材 & 理论学习基础，并为知识图谱和教学内容提供教材依据 \\
    专家共识与临床规范 & 实操规范，为诊疗选择、适用条件和风险控制提供权威依据 \\
    《默沙东诊疗手册》 & 补充各类疾病、药物、症状鉴别和跨学科的医学知识 \\
    口腔执业医师考试相关题库 & 课后练习和测试，支持薄弱点和知识盲区的识别 \\
\end{longtable}
}

数据输入先按照资料形态完成基础预处理，保留表格、版面、文字与层级结构，并统一整理文档、页面、片段和段落位置。文本内容按照引文范围和主题边界合并为语义单元，教材知识图谱、医学文本和题库随后建立双向联系，再生成统一检索语料以及稀疏、向量和图关系索引。标识与引用校验、重复内容标记、链接验证和质量检查随各处理步骤执行，整体流程如图~\ref{fig:data-processing-pipeline} 所示。

\begin{figure}[H]
  \centering
  \includegraphics[trim=150 8 150 8,clip,width=\linewidth]{kb-data-process.pdf}
  \caption{OralSage 多源数据处理流程}
  \label{fig:data-processing-pipeline}
\end{figure}

\subsection{口腔医学教材}

项目收集整合了 14 本口腔医学权威教材，这些教材多选自“人卫十三五”规划系列，涵盖口腔医学的临床诊疗、生物机理、材料科学及全流程治疗方案等核心领域，侧重于基础理论的深度阐释和学科体系的完整性，构成了系统的“学术基础层”，为学习内容提供完整的学科体系、稳定的术语规范和教学级知识依据。教材具体说明如表~\ref{tab:textbook-inputs}所示。
{\small
\renewcommand{\arraystretch}{1.08}
\begin{longtable}{p{0.27\linewidth}p{0.10\linewidth}p{0.49\linewidth}}
  \caption{口腔医学教材输入清单}\label{tab:textbook-inputs}\\
  \toprule
  教材 & 版本 & 主要知识范围 \\
  \midrule
  \endfirsthead
  \multicolumn{3}{c}{续表~\thetable}\\
  \toprule
  教材 & 版本 & 主要知识范围 \\
  \midrule
  \endhead
  \bottomrule
  \endfoot
  《儿童口腔医学》 & 第5版 & 儿童口腔生长发育、龋病、牙髓根尖周病与牙外伤 \\
  《口腔材料学》 & 第6版 & 材料组成、理化性能、生物学性能与临床选择 \\
  《口腔颌面外科学》 & 第8版 & 颌面部疾病、创伤、肿瘤、诊断与外科操作 \\
  《口腔颌面医学影像诊断学》 & 第7版 & 影像检查原理、正常影像表现与疾病影像诊断 \\
  《口腔临床药物学》 & 第5版 & 口腔常用药物、药理作用、适应证与用药风险 \\
  《口腔黏膜病学》 & 第5版 & 口腔黏膜疾病的病因、表现、诊断与治疗 \\
  《口腔生物学》 & 第5版 & 口腔组织的细胞、分子、微生态与生物学机制 \\
  《口腔修复学》 & 第8版 & 牙体、牙列及颌面缺损的修复设计与临床流程 \\
  《口腔预防医学》 & 第7版 & 口腔流行病学、疾病预防、健康教育与社区保健 \\
  《口腔正畸学》 & 第7版 & 错颌畸形的病因、诊断、矫治设计与生物力学 \\
  《口腔种植学》 & 第1版 & 种植生物学、手术、修复设计与并发症管理 \\
  《口腔组织病理学》 & 第8版 & 口腔组织发育、正常结构及疾病病理变化 \\
  《牙体牙髓病学》 & 第5版 & 龋病、牙体硬组织病、牙髓病与根尖周病 \\
  《牙周病学》 & 第5版 & 牙周组织、牙周疾病的诊断、治疗与维护 \\
\end{longtable}
}

由于口腔医学各学科之间具有较强的知识依赖，真实的口腔医学课程学习中，学生往往需要建立起清晰的知识架构，同步理解正常组织结构、生物学机制、病理变化和临床表现。因此，本项目在处理教材资料时保持从基础理论到临床应用的连续知识脉络，将疾病分类、专业术语和诊疗原则经过系统编排，通过连续的章节论述，保留概念产生的背景与前后联系，避免零散知识片段带来的割裂感。

由于教材原始资料多为图片扫描版PDF，无法直接进行文本检索，为了将教材中的学科结构和医学内容转换为可检索定位的教学证据，本项目具体处理方法如下。

\noindent{（1）图像增强与 OCR 识别}
\nopagebreak[4]

教材 PDF 的原始扫描件存在分辨率不一、噪点较多等问题。首先对扫描页进行图像增强和高清压缩，随后使用开源 PaddleOCR 工具批量提取文字，并将识别结果统一转换为 UTF-8 文本。针对医学文档中常见的复杂排版和专有名词，对识别结果进行人工抽检和校对，将非图片内容转化为可编辑的文字模态信息。

\noindent{（2）保留章节结构与来源位置}

教材知识依赖章、节和局部标题形成连续语境，仅保存正文会丢失知识所属位置。因此，在对文字进行提取时保留章标题、节标题、局部标题、图题、表题和页码，并为教材、页面和局部文本片段建立逐级位置关系。

\noindent{（3）语义分块与索引复用}

由于 OCR 文本中的换行可能来自真实段落，也可能来自页面排版。本项目在处理时根据空行、章节编号、短标题、括号标题、编号条目和图表提示判断行的类型，同时清理重复页眉、页码和低信息噪声。对属于同一语句的连续行进行重新连接，章标题、节标题和局部标题则单独保留，并按出现顺序形成标题路径，使后续正文能够继承距离最近的有效标题。此外，按照语义构建分块，优先在新的章、节或主要标题处结束当前语义块，在必要时调整块的长度，确保按照知识逻辑进行组织，为后续知识图谱节点、检索结果和生成内容保留可定位的教材内容。索引构建后通过了召回验证，重点检查了专业术语和疾病分类的检索准确率，确保知识库在实际使用中具备基本可靠性。

\subsection{专家共识与临床规范}

本项目收集了 30 项与临床实践紧密相关的口腔医学专家共识、诊疗指南和操作规范，其中包括《创伤性周围性面瘫外科治疗专家共识》《单侧唇裂诊疗指南》《牙周龈下刮治和根面平整术操作规范》和《口腔种植体周病管理专家共识》等资料（详见附件）。涉及学科覆盖口腔颌面外科、牙体牙髓、牙周、修复、正畸、种植、口腔黏膜、儿童口腔和急诊处置。

专家共识代表了行业内的“金标准”，具有表述规范、逻辑严密、指导性强的特点。这部分数据作为临床规范，为适应证判断、治疗方案制定和操作风险控制提供权威依据。专家共识的补充，可以进一步回答特定条件下应选择何种方案、何时不宜实施以及操作中需要控制哪些风险。在处理临床应用题、诊疗原则和操作规范问题时，以临床规范为约束，能够减少智能体仅凭一般医学知识推演临床结论的风险，增强诊疗内容的严谨性。

专家共识和诊疗指南多以 PDF 格式输入，存在大量页眉页脚、分页残留和多栏分布文字，为了最大限度保留专家共识和数据指南材料的医学严谨性并适配知识库检索，制定以下处理策略。

\noindent{（1）文本清理与格式矫正}

处理过程中需要清理重复标点与异常空白，并按照栏区恢复双栏或多栏页面的阅读顺序。对于跨页但语义连续的段落，需要重新连接前后内容，使推荐意见及其适用条件能够保持完整。

\noindent{（2）临床条款结构保留}

专家共识在引用时需要与适用条件和操作要求共同出现，缺少上下文容易改变临床含义，且该类资料一般以条款式或分层级小节形式呈现，章节之间逻辑紧密，但每一段内部内容相对独立，较适合自动切分。因此，在处理过程中保留“推荐意见”“注意事项”“适应证”“禁忌证”和“操作步骤”等标题及其条款编号，并将标题与后续正文共同组织。经过处理后，治疗建议、适用范围和风险提示仍然保持明确对应关系。

\noindent{（3）表格语义重组}

诊疗指南中的表格将适用条件、操作要求和评价结果分置在不同列，逐单元格提取会破坏同一条临床信息的联系。因此，处理时按照数据行或完整逻辑单元转写表格，使同一项诊疗条件和结论进入连续文本，并保留表题及其所在章节，从而为后续分块提供完整语义。

\subsection{默沙东诊疗手册}

《默沙东诊疗手册》作为全球连续出版时间最长的综合性医学手册之一，自1899年首次发行以来，已成为世界上使用最广泛的医学信息资源之一。其核心特征包括内容全面性、权威可靠性、版本多样性、免费获取性和更新及时性，为医疗专业人士和普通公众提供科学准确的医学信息支持，在本项目中作为通用医学知识的重要来源。

口腔症状可能与感染、免疫、代谢和肿瘤等全身性问题相关，患者用药也会影响口腔表现与治疗选择。手册中的疾病百科、症状解析、检查方法、治疗原则和药物信息，能够补充基础病理知识缺失，使系统能够把口腔专科问题放入更完整的医学背景中理解。基础疾病查询、症状鉴别和药物咨询可从中获取跨学科线索，教材中的口腔知识也能由此联系炎症反应、组织损伤和全身疾病机制。

为了将手册内容适配到本项目知识库检索系统中，制定以下数据处理方法。

\noindent{（1）数据采集与格式处理}

使用爬虫技术将默沙东诊疗手册的公开内容完整爬取，统一文字编码与空白格式，并清理导航、模板和其他网页噪声，初步整理为JSONL格式以便存储。

\noindent{（2）保留疾病条目层级}

疾病条目中的概述、症状、检查、治疗和药物信息具有明确的从属关系，脱离标题后容易失去医学语境。因此，处理过程中按照网页标题层级组织正文，并把疾病名称和局部主题补充到对应文本中，使每段内容仍能确定所属疾病及其医学主题，为全身疾病口腔表现、药物作用、症状鉴别和跨学科概念解释提供补充依据。


\subsection{口腔执业医师考试相关题库}

题库面向口腔医学执业医师考试相关科目，共收集 3,161 道选择题，覆盖 12 个学科。题量分布见表~\ref{tab:question-inputs}，其中临床医学题目用于补充口腔疾病相关的全身医学基础。

\begin{table}[H]
  \centering
  \caption{执业医师考试相关题库的学科分布}
  \label{tab:question-inputs}
  \begin{tabularx}{\linewidth}{X r X r}
    \toprule
    学科 & 题量 & 学科 & 题量 \\
    \midrule
    牙体牙髓病学 & 651 & 口腔颌面外科学 & 549 \\
    口腔修复学 & 436 & 牙周病学 & 435 \\
    临床医学 & 410 & 口腔黏膜病学 & 204 \\
    儿童口腔医学 & 120 & 口腔组织病理学 & 106 \\
    口腔预防医学 & 104 & 口腔解剖生理学 & 72 \\
    口腔材料学 & 44 & 口腔颌面医学影像诊断学 & 30 \\
    \midrule
    合计 & \multicolumn{3}{r}{3,161} \\
    \bottomrule
  \end{tabularx}
\end{table}

题干将病史、症状、检查结果和治疗条件组织为具体的临床情境，要求学习者能够识别影响诊断或治疗选择的关键信息，并运用相应知识完成判断。该类题目精心设计了大量混淆项，不同选项之间往往只有疾病阶段、适用条件或处置原则上的细微差别，能够有效暴露出学习者的真实学习情况。学习者在答题过程中如果进行了错误的选择，在一定程度上能体现其该类知识点上的薄弱之处，系统会予以记录并更新学习画像，为学习者提供针对性反馈。

为了从题库 PDF 中得到题目结构完整、答案对应准确的标准记录，本项目具体处理方法如下。

\noindent{（1）恢复 PDF 分栏顺序}
\nopagebreak[4]

题库 PDF 采用双栏排版，直接提取容易把相邻题目的题号、题干和选项交叉拼接。程序在识别文字前按照栏区恢复阅读顺序，并检查题号与选项是否连续，使每道题目的内容保持在同一记录中。

\noindent{（2）对齐题目、答案与解析}

项目所收集到的资料中，题目和解析大多分开编排，答案汇总表与逐题解析也采用不同格式。在处理时需要识别题目页、答案汇总表和解析段落，再按照科目、章节和章节内题号建立对应关系。经过对齐后，每道题都能够同时关联题干、选项、答案和解析，并保留其在 PDF 中的题目块、答案块和解析块位置。


\noindent{（3）建立题目与知识点关联}

首先在同名科目和对应章节内匹配知识图谱中的核心知识、支撑知识和概念，再根据节点归属关系连接到正式核心知识点。处理后形成 3,161 个题目对象和 6,287 条题目到知识点的关系，使题目能够沿知识节点访问相应教材证据，并为练习推荐、错因分析和掌握度更新提供结构化依据。


\section{知识图谱构建}

完成教材文本整理和章节位置恢复后，项目使用 OralPath KG-Forge 构建口腔医学知识图谱。该流程借鉴编译系统的处理方式，将版本化教材与指南作为输入，将书、章、节组成的文档结构树作为章节映射，将受控医学 Schema 作为类型规则，并以片段标识、双页码和连续引文记录知识来源。模型抽取结果先写入 Evidence-to-Graph Intermediate Representation（EGIR），经过医学审议与确定性编译后，再生成口腔医学知识图谱节点。

\subsection{自动化持续构建}

构建过程中，系统能够实现新教材接入、教材版本更新和指定章节重建。Corpus Watcher 为文档、章节和局部片段计算内容指纹，并与上一次构建清单比较，从而确定需要处理的章节范围。文档随后解析为书、章、节组成的结构树，每个层级保留文档版本、PDF 页序、印刷页码和局部片段位置，章节范围按照同级标题边界确定，后续抽取能够始终在明确的教材范围内进行。

Incremental Build Planner 根据内容指纹和来源依赖关系生成持久化任务 DAG，将构建批次拆分为章节抽取、章节检查、跨章节归并和制品发布等任务。彼此独立的章节由 Worker Pool 并行处理，来源内容、Schema 版本和抽取配置共同形成任务幂等键。内容未发生变化的章节复用既有结果，发生变化的章节及其关联关系进入重建流程，使多本教材能够按照最小必要范围持续更新。

\subsection{Schema 引导抽取与 EGIR 生成}

提取章节过程中，首先识别普通医学、指南与防护、绪论与概述等不同内容模式，再选择对应的节点类型和关系集合。普通医学章节围绕疾病、病因、病原体、风险因素、症状、体征、影像表现、病理变化、诊断方法、治疗、并发症和解剖结构组织知识。放射防护、操作安全和规则型章节增加防护原则、防护措施、剂量限制、设备要求、操作规则和辐射风险等类型，合理处理规范化内容。

模型根据章节文本生成 EGIR，其中保存候选节点、原子医学陈述、补充概念、来源指针和抽取记录。每条候选陈述在生成时即关联连续引文、片段标识和页码位置，使实体、关系与教材内容从抽取阶段开始保持对应。书、章、节继续承担来源结构的组织功能，正式医学实体从核心知识层开始，章节标题不会混入疾病、治疗或诊断节点。

\subsection{节点与关系编译}

知识对象按照核心知识、支撑知识和补充概念分层组织，无法进入图谱的候选项同时保留处理原因。各层的编译要求见表~\ref{tab:kg-object-contract}。

\begin{table}[H]
  \centering
  \caption{知识图谱对象分层与编译要求}
  \label{tab:kg-object-contract}
  \begin{tabularx}{\linewidth}{>{\raggedright\arraybackslash}p{0.18\linewidth}p{0.40\linewidth}X}
    \toprule
    知识对象 & 医学内容 & 编译要求 \\
    \midrule
    核心知识 & 章节或节次重点讨论的疾病、解剖结构及医学主题 & 使用规范名称和受控类型，每个核心知识至少参与两条正式关系。 \\
    支撑知识 & 病因、症状、体征、影像、病理、诊断、治疗和并发症等关系对象 & 通过正式关系连接核心知识，每个支撑知识至少参与一条正式关系。 \\
    补充概念 & 来源明确且适合辅助解释，但难以使用受控谓词完整表达的内容 & 使用 \texttt{RELATE} 连接一至三个相关节点，并记录采用补充连接的原因。 \\
    舍弃候选 & 引文不完整、内容重复、缺少合理连接或图谱信息量较低的候选项 & 不写入正式节点与关系，同时保留候选内容、处理原因和对应引文。 \\
    \bottomrule
  \end{tabularx}
\end{table}

正式关系采用受控的主语、谓词和宾语结构。普通医学关系覆盖病因、病原体、风险因素、症状、临床体征、影像表现、病理特征、并发症、诊断方法、治疗方法、发生部位、鉴别对象和疾病亚型。指南与防护关系补充适用对象、前置要求、防护原则、防护措施、剂量限制、设备要求和风险降低等语义。编译程序同时检查谓词方向以及两端节点类型，使疾病与病因、疾病与诊断方法、操作规则与适用对象等关系保持一致的表达方式。

\subsection{医学审议与符号编译}

EGIR 依次经过医学抽取、类型复核、来源核验、图结构检查和综合裁决。医学抽取单元将连续文本拆为原子医学陈述并补全候选节点，类型复核单元检查节点类别、关系谓词和连接方向，来源核验单元判断连续引文能否完整支持一条主谓宾关系，图结构检查单元识别重复节点、自环、悬空连接和概念挂接异常，综合裁决单元据此确定正式关系、补充概念或舍弃候选。

完成审议的 EGIR 随后进入确定性编译步骤。程序依次规范实体名称和别名，检查节点与关系类型，核对片段、页码和连续引文，合并重复节点与完全相同的关系，并检查核心知识度数、支撑知识连接和补充概念挂接。可修正的名称、引用和连接问题进入一次有界修复，修正结果重新接受同一组检查。仍未满足结构要求的候选连同诊断信息进入隔离队列，其余章节可以继续完成构建与发布。

质量检查由 Schema 符合度、医学类型符合度、来源覆盖、引文支持、图结构完整性、概念连接完整性和构建可复现性共同组成。检查过程采用明确规则与结构化诊断，节点数量不会单独决定图谱质量。每条正式关系均保留片段标识、PDF 页序、印刷页码、连续引文和关系说明，使图谱中的医学陈述能够定位到对应教材内容。

\subsection{跨章节归并与交付}

各章节完成编译后，Resolution Job 根据规范名称、别名、节点类型和章节语境识别跨章节同名实体，并生成统一映射。章节内部的局部引用保持不变，单章节图谱可以独立读取，跨章节检索也能够识别指向同一医学实体的多个局部节点。该处理兼顾章节制品的独立性与全书范围内的术语一致性。

Publish Job 将章节 JSON、文档级索引和质量报告作为同一发布单元，解析检查完成后再登记正式版本。每次构建记录来源文档与版本、内容指纹、Schema 版本、流水线配置、制品摘要、质量报告和依赖快照。相同输入与配置对应相同的构建键，新版本以独立制品保存，使图谱生成过程具备清晰的版本链和来源记录\cite{w3c-prov-dm}。

系统同时维护局部片段到节点、关系和章节制品的依赖关系。教材版本或 OCR 文本发生变化时，Dependency Tracker 标记受影响的图谱对象，并由构建规划器重新处理相关章节与关系。未受影响的章节继续复用既有制品，从而完成按来源变化范围执行的增量更新。发布后的图谱节点、正式关系、补充概念、章节归属和来源位置随后进入 6.3 节的统一入库与索引构建流程。

\section{知识库入库与索引组织}

完成多源资料预处理和知识图谱构建后，将教材语义文本单元、临床规范文本、诊疗手册文本、知识图谱节点和题库题目进行知识库入库，核心工作包括检索对象组织、多路索引构建和知识关系连接，为后续 KAG 求解器提供候选召回、关系扩展和上下文基础。

\subsection{检索对象的统一组织}

多源资料入库后形成教材语义文本单元、临床规范文本、诊疗手册文本、知识图谱节点和题库题目五类检索对象。文本单元提供教材、规范和手册中的医学原文内容，知识图谱节点提供口腔医学概念、章节归属和关系结构，题库题目提供考查场景、选项解析和学习反馈依据。经过统一组织表示后，检索器能够在同一流程中处理医学正文、结构化知识和题库材料，并将不同来源的候选内容组织为面向 KAG 求解器的一致输入。

\subsection{多路索引的构建机制}

口腔医学查询包含明确术语、病例化描述、同义表达和关系推理需求。本项目围绕入库对象构建术语别名索引、稀疏倒排索引、稠密向量索引和图关系索引。四类索引分别承担名称匹配、词项匹配、语义匹配和关系扩展功能，其形式化表示为：

\begin{equation}
  \label{eq:multi-index-set}
  \mathcal{I}=\{\mathcal{I}_{\mathrm{alias}},\mathcal{I}_{\mathrm{sparse}},\mathcal{I}_{\mathrm{dense}},\mathcal{I}_{\mathrm{graph}}\}
\end{equation}

（1）术语别名索引。Alias Index 由知识节点的规范名称、常用别名、英文术语和专业缩写构成，采用字符串规范化与别名表匹配方式进行索引。其目的在于处理疾病名、操作名、解剖结构名和专业缩写等明确查询，并将查询直接映射到候选知识节点，为后续图扩展提供初始种子。

（2）稀疏倒排索引。BM25 Sparse Index 采用中文分词、英文数字识别和中文双字组合构建倒排结构，并按照 BM25 评分衡量查询词项与候选文本之间的匹配程度。该索引用于教材原文、题干关键词、数值条件和医学术语召回，能够保留精确词项匹配对医学检索的约束作用。

（3）稠密向量索引。BGE-M3 Dense Index 使用 BGE-M3 模型\cite{chen2024m3embedding}对入库对象进行语义编码，并通过归一化向量计算查询与候选对象之间的语义相似度。FAISS 进行相似向量检索功能，BGE-M3 负责文本表示生成。该索引用于同义表达、口语化问题和病例化描述召回，补充稀疏匹配难以覆盖的语义相关内容。

（4）图关系索引。Graph Index 保存知识节点、医学文本单元和题库题目之间的邻接关系，覆盖 KG adjacency、KG-Chunk links 和 Question-KG links 等关系。该索引用于从候选知识节点扩展到相关概念、医学原文和练习题目，并支持 KG path 输出、Graph-PPR 候选传播和来源定位。

\subsection{知识关系与来源关联}

入库阶段将知识图谱、医学文本和题库题目连接为能够遍历的关系结构。教材知识节点与医学文本单元之间建立知识到文本的关联，使知识点能够返回对应的原文材料。题库题目与知识节点之间建立考查关系，使练习推荐、错因分析和掌握度更新能够回到具体知识点。

医学文本单元继续保留资料、页码、span 和标题路径等来源信息。检索命中文本单元时，系统能够沿关联知识标识返回相关知识点。术语别名索引命中知识点时，系统能够沿关系结构获得相关医学文本和题目材料。图关系索引将这些连接组织为在线能够遍历的邻接结构，为 Graph-PPR、KG path 输出和引用定位提供基础。

\subsection{入库结果}

入库完成后，知识库形成了统一组织医学文本、结构知识、题库材料和来源位置的知识对象空间。专家共识和《默沙东诊疗手册》作为独立来源层保留各自的资料属性，用于补充口腔专科问题中的临床规范和通用医学参考。各类对象在入库阶段完成标识统一、关系连接和索引准备，具体在线召回、扩展与排序过程在 6.6 节中说明。

面向上层智能体，入库结果能够提供医学文本、相关知识点、关系路径和来源位置。不同智能体通过统一检索接口读取与当前任务相关的对象，具体调用方式在 6.6 节的 KAGSolver 与智能体调用部分说明。

\section{KAG 知识库架构}
本项目中的知识库需要同时支持课程规划、多模态资源生成、教学讲解、智能答疑和学习测评等任务。基于文本向量的检索方式难以同时表达教材章节关系、诊疗知识联系和题目考查关系，也难以让不同学习任务使用同一组知识对象。为此，项目采用基于 KAG 的知识库架构，沿用知识图谱与文本互索引、多路检索和图关系扩展等基础设计，并结合口腔医学资料特点补充教材章节组织、题目关联、来源属性和多智能体调用接口。

\subsection{KAG 基础架构与项目适配}

KAG 面向专业知识服务，将知识图谱的结构表达与向量检索的语义匹配能力结合起来，并通过适合大模型使用的知识表示、知识图谱与文本互索引、混合推理以及知识对齐，缩小文本相似度与知识推理相关性之间的差距\cite{liang2024kag}。本项目采用其中的结构化知识表示、双向索引和混合检索思路，以前期构建的口腔医学教材知识图谱组织学科范围与知识关系，通过结构化存储的医学文本保留专业表述，再通过多路检索与关系扩展为上层智能体整理学习上下文。

项目知识库分为知识表示、知识互索引、混合检索和知识服务四个部分。知识表示负责保存教材层级、知识点及医学关系；知识互索引负责连接知识节点、语义文本单元和考试题目；混合检索负责从不同匹配角度取得相关对象；知识服务则将结果整理为上层智能体能够按任务读取的统一上下文。针对口腔医学学习场景，项目进一步补充资料层级、题目到知识点的考查联系和多智能体调用接口，具体对应关系见表~\ref{tab:kag-adaptation}。

\begin{table}[H]
  \centering
  \caption{KAG 基础能力与 OralSage 场景适配}
  \label{tab:kag-adaptation}
  \begin{tabularx}{\linewidth}{p{0.22\linewidth}p{0.38\linewidth}X}
    \toprule
    KAG 基础能力 & OralSage 场景适配 & 支撑任务 \\
    \midrule
    结构化知识表示 & 以教材学科、章节和知识点层级组织口腔医学概念及其关系 & 学习顺序规划与知识范围控制 \\
    知识与文本互索引 & 建立知识点、医学内容和资料位置之间的双向联系 & 讲义生成、教学讲解与来源定位 \\
    混合检索与图扩展 & 联合术语、词项、语义向量和知识关系获取候选内容 & 病例化提问、同义表达与关联知识检索 \\
    知识服务输出 & 区分教材、临床规范、外部参考和题库的资料属性，并组合医学内容、知识节点、关系路径、引用位置和检索轨迹 & 多智能体调用、医学内容组织与学习评价反馈 \\
    \bottomrule
  \end{tabularx}
\end{table}

\subsection{知识库总体架构}

图~\ref{fig:kb-architecture} 将 KAG 知识库划分为知识库构建流、异构知识空间、KAG Solver 检索链路、来源与质量控制以及多智能体调用五个相互连接的区域。资料沿左侧构建流进入中央知识空间，查询从右侧检索链路获得经过组织的回答上下文，顶部智能体通过统一接口调用这些结果，底部的来源与质量控制为构建和检索过程提供限制条件。

\begin{figure}[H]
  \centering
  \includegraphics[width=\linewidth]{kb-architecture.pdf}
  \caption{OralSage KAG 知识库总体架构}
  \label{fig:kb-architecture}
\end{figure}

图中左侧构建流对应前述多源资料处理、知识图谱构建和统一入库过程，中央异构知识空间负责保存知识节点、关系、文本单元、来源片段和题目对象。右侧 KAG Solver 将这些对象组织为可检索的知识服务，其召回、图扩展、融合排序和结果补全过程在 6.6 节中展开说明。

底部来源与质量控制贯穿对象构建和在线检索，负责保持知识对象、文本内容与资料位置之间的对应关系；顶部多智能体调用区通过统一接口获取已经组织的检索上下文。由此，本节侧重说明各区域在总体架构中的位置与协作关系，具体算子封装见 6.5 节，检索约束和内容审查见 6.7 节。

\subsection{六层知识服务组织}

最终，本项目知识库组织为六个层次，如图~\ref{fig:kb-layered-architecture} 所示，依次完成资料定位、语义组织、图谱关联和索引生成。

\begin{figure}[H]
  \centering
  \includegraphics[width=\linewidth]{kb-layered-architecture.pdf}
  \caption{OralSage KAG 知识库分层结构}
  \label{fig:kb-layered-architecture}
\end{figure}


\section{自定义算子设计}

\subsection{算子设计规范}

OralSage 将知识库构建与在线检索拆分为若干具有明确输入和输出的算子。这里的自定义主要体现为对资料对象、知识图谱对象、索引资产和智能体接口进行统一编排，使 KAG 的基础组件能够接入本项目的数据处理流程。各算子沿用统一的对象标识和元数据格式，输入中记录资料类型、来源位置和质量标记，输出中保留对象之间的关联关系。

每个算子采用独立名称标识，并说明所属处理阶段、输入对象和输出对象。离线算子按照确定的顺序生成文本单元、知识对象和索引资产，并将运行结果写入构建清单，便于检查数据是否完整。在线算子以查询和已生成的索引为输入，返回候选内容、知识关系、资料位置和检索记录。

\subsection{数据与索引构建算子}

数据与索引构建算子按照资料规范化、语义单元生成、知识对象关联、索引校验的顺序执行。前两步将不同格式的来源资料转换为可检索的文本单元，第三步连接教材知识图谱、文本和题库，最后一步生成多路索引并检查各类对象之间的对应关系。

\begin{table}[H]
  \centering
  \caption{知识库数据与索引构建算子}
  \label{tab:kb-build-operators}
  \begin{tabularx}{\linewidth}{>{\raggedright\arraybackslash}p{0.24\linewidth}>{\raggedright\arraybackslash}p{0.20\linewidth}X}
    \toprule
    算子名称 & 处理阶段 & 输入与输出 \\ 
    \midrule
    Source Normalizer & 资料规范化 & {输入：} PDF、网页文本、OCR 文本、题库记录和来源信息。\newline {输出：} 文档记录、清理文本、标题路径、资料位置和质量标记。 \\ 
    Span Semantic Chunker & 语义单元生成 & {输入：} 段落单元、标题路径、引文范围和长度参数。\newline {输出：} 按主题边界组织的语义文本单元、来源片段和局部位置。 \\ 
    Canonical Object Builder & 知识对象规范化 & {输入：} 教材知识图谱草稿、题库匹配结果和资料属性。\newline {输出：} 具有统一标识和章节归属的知识节点、知识关系、概念关系、别名和题目对象。 \\ 
    KG--Chunk Linker\newline Question--KG Linker & 对象关联 & {输入：} 知识节点、文本单元、来源片段和题目对象。\newline {输出：} 知识到文本、文本到来源、题目到知识点的链接。 \\ 
    Sparse BM25 Builder\newline Dense BGE-M3 Builder\newline Graph Index Builder & 多路索引生成 & {输入：} 统一检索语料、BGE-M3 模型、知识关系和别名。\newline {输出：} BM25、向量、FAISS、图邻接和别名索引。 \\ 
    Index Asset Validator & 资产校验 & {输入：} 正式对象、索引和构建清单。\newline {输出：} 对象标识、链接目标、文档顺序和向量顺序的校验结果。 \\ 
    \bottomrule
  \end{tabularx}
\end{table}


\subsection{检索与结果组织算子}

在线检索被封装为候选召回、图扩展、候选筛选、融合排序、内容补全和上下文组织等具有明确输入输出的算子，如表~\ref{tab:kb-query-operators}所示。本节侧重说明各算子的职责边界，具体检索公式、参数和执行过程在 6.6 节中展开。

\begin{table}[H]
  \centering
  \caption{KAGSolver 在线检索算子}
  \label{tab:kb-query-operators}
  \begin{tabularx}{\linewidth}{>{\raggedright\arraybackslash}p{0.24\linewidth}>{\raggedright\arraybackslash}p{0.20\linewidth}X}
    \toprule
    算子名称 & 处理阶段 & 输入与输出 \\ 
    \midrule
    Alias Seed\newline Sparse BM25\newline Dense BGE-M3 & 候选召回 & {输入：} 查询、别名索引、BM25 索引和 FAISS 索引。\newline {输出：} 术语命中、关键词重合和近义表达对应的候选集合。 \\ 
    Graph Seed Collector\newline Graph-PPR Expansion & 图种子汇聚与扩展 & {输入：} 多路候选和图邻接。\newline {输出：} 知识节点种子、图候选、局部关联对象和关系路径。 \\ 
    Hybrid Union\newline Quality Gate & 候选合并与筛选 & {输入：} 多路候选、资料属性和质量标记。\newline {输出：} 按统一对象标识合并并完成质量筛选的候选集合。 \\ 
    RRF Fusion & 融合排序 & {输入：} 各通道排名和候选顺序。\newline {输出：} RRF 排序结果与融合分数。 \\ 
    Evidence Hydration\newline Citation Builder\newline KG Path Builder & 内容与路径补全 & {输入：} 排序后的对象标识、知识关系和来源位置。\newline {输出：} 智能体可读的完整文本、知识节点、关系路径和引用位置。 \\ 
    Agent Payload Builder & 上下文组织 & {输入：} 补全后的对象、排序信息和检索记录。\newline {输出：} 面向 Planner、Resource Factory、Tutor、QA 和 Assessment 的统一检索上下文与检索轨迹。 \\ 
    \bottomrule
  \end{tabularx}
\end{table}

算子之间通过统一的对象标识传递结果，保留参与的召回通道、候选数量、图扩展结果、排序方式和资料位置，便于分析一次检索由哪些步骤产生，集中处理检索和结果组织。

\section{知识库检索}

\begin{figure}[H]
  \centering
  \includegraphics[width=\linewidth]{kb-search-process.pdf}
  \caption{OralSage 知识库检索流程}
  \label{fig:kb-search-process}
\end{figure}

知识库检索按照图~\ref{fig:kb-search-process} 所示的四个阶段执行。查询经过术语规范化后进入 Alias、BM25 和 BGE-M3 三路召回，各通道返回的候选映射到知识节点，并由 Graph-PPR 沿局部关系扩展相关节点、医学文本和题目。多路结果按照统一对象标识合并，经过质量筛选和 RRF 排序后补充完整医学文本、知识节点、资料位置、关系路径与引用信息，最终形成供 Planner、Resource Factory、Tutor、QA 和 Assessment 调用的统一检索上下文。


\subsection{多路候选召回}

KAGSolver 并行执行术语别名、BM25 和 BGE-M3 三路召回。术语别名通道直接匹配知识节点名称、别名和缩写，适合识别规范疾病名称、检查术语和治疗名称。BM25 通道对中英文内容进行词项切分，并保留医学术语、数字、英文缩写和教材表述之间的精确重合，适合处理定义查询和条件明确的问题。BGE-M3 通道对查询与检索文本进行语义编码，用于发现表达方式不同但医学含义接近的内容。

BM25 使用平滑后的逆文档频率衡量词项区分度。设语料库文档总数为 \(N\)，包含词项 \(t\) 的文档数为 \(df(t)\)，则：

\begin{equation}
  \label{eq:bm25-idf}
  \operatorname{idf}(t)=\log\left(\frac{N-df(t)+0.5}{df(t)+0.5}+1\right)
\end{equation}

对于查询 \(q\) 和文档 \(d\)，BM25 文档得分表示为：

\begin{equation}
  \label{eq:bm25-score}
  \operatorname{Score}_{\mathrm{BM25}}(q,d)
  =\sum_{t\in q}qtf(t)\operatorname{idf}(t)
  \frac{tf(t,d)(k_1+1)}
  {tf(t,d)+k_1\left(1-b+b\frac{|d|}{\operatorname{avgdl}}\right)}
\end{equation}

式中，\(tf(t,d)\) 和 \(qtf(t)\) 分别表示词项在文档和查询中的出现次数，\(|d|\) 表示文档长度，\(\operatorname{avgdl}\) 表示平均文档长度。系统设置 \(k_1=1.5\)、\(b=0.75\)。

BGE-M3 的最大输入长度为 512 个词元，编码后得到 1,024 维向量。设 \(E(x)\) 为文本 \(x\) 的编码结果，向量归一化与语义相关度计算为：

\begin{equation}
  \label{eq:dense-score}
  \hat{E}(x)=\frac{E(x)}{\lVert E(x)\rVert_2},
  \qquad
  \operatorname{Score}_{\mathrm{dense}}(q,d)=\hat{E}(q)^{\top}\hat{E}(d)
\end{equation}

归一化后的内积等价于余弦相似度，FAISS 据此返回语义距离较近的文本。三路召回结果保留各自的匹配方式和排名信息，为图关系扩展与融合排序提供输入。

\subsection{图种子汇聚与 Graph-PPR}

图种子汇聚模块将三路召回结果转换为知识节点。术语别名命中的节点直接加入种子集合，BM25 或 BGE-M3 命中知识节点时使用对应节点标识，命中文本单元或题目时则沿对象链接取得相关知识点。由此，名称匹配、词项匹配和语义匹配均能进入同一图扩展过程。

Graph-PPR 采用 Personalized PageRank 的局部随机游走思想\cite{page1999pagerank}。设 \(S\) 为查询种子分布，\(P\) 为图转移矩阵，\(\alpha\) 为重启概率，则迭代过程为：

\begin{equation}
  \label{eq:graph-ppr}
  r_{t+1}=(1-\alpha)P^{\top}r_t+\alpha S
\end{equation}

系统设置 \(\alpha=0.18\)，执行 18 次迭代，并从节点关系、概念关系、知识到文本关系和题目到知识点关系中取得局部候选。关系路径保留查询种子到结果对象的三跳以内连接，用于说明内容之间的知识联系。

\subsection{候选融合与 RRF 排序}

术语别名、BM25、BGE-M3 和 Graph-PPR 返回的候选按照统一对象标识合并，相同对象只保留一条记录，同时汇总其匹配通道和图关系得分。候选筛选读取资料属性与质量标记，降低重复文本、模板内容和参考文献密集片段的排序优先级。完成合并与筛选后，系统使用倒数排名融合方法计算综合顺序\cite{cormack2009rrf}：

\begin{equation}
  \label{eq:rrf}
  \operatorname{Score}_{\mathrm{RRF}}(d)=\sum_{r\in R_d}\frac{1}{k+\operatorname{rank}_r(d)}
\end{equation}

其中，\(R_d\) 表示对象 \(d\) 进入的排名通道集合，\(\operatorname{rank}_r(d)\) 表示该对象在通道 \(r\) 中的名次，平滑常数 \(k\) 取 60。一个对象由多个通道共同召回时，其综合得分会相应提高，从而兼顾精确术语命中、语义相关性和图关系关联。

\subsection{检索结果组织}

排序完成后，KAGSolver 根据对象标识读取完整内容，并组织为医学文本、相关知识节点、关系路径、引用位置和检索轨迹。医学文本保留资料类型、标题路径和来源片段，引用位置记录资料名称、页码范围及局部片段，关系路径记录查询种子到结果对象之间的连接。检索轨迹则汇总参与的召回通道、候选数量、图种子数量和融合分数。

这些内容采用统一结构返回，使上层调用既能取得直接用于生成的医学文本，也能读取与其关联的知识点和资料位置。结果组织阶段按照返回数量控制上下文规模，并保留排序分数和匹配方式，便于不同学习任务选择合适内容。

\subsection{KAGSolver 与智能体调用}

后端将 KAGSolver 作为进程内单例使用，首次调用加载检索语料、BM25、FAISS、图邻接、别名索引和 BGE-M3 查询编码器，后续请求复用已加载的索引与模型。统一检索接口将 KAG 返回结果映射为 \texttt{EvidenceChunk} 对象，使各智能体使用一致的文本、来源和分数格式。

Planner 在课程规划前检索主题相关知识，并结合章节归属和关系路径确定学习范围。Resource Factory 根据讲义小节组织定向查询，Tutor 与 QA 调用医学文本和引用位置完成讲解与答疑，Assessment 通过题目到知识点的关系分析错因并更新掌握情况。各智能体可以改变查询内容和返回数量，同时复用同一套检索链路与结果结构。

\section{防幻觉与内容审查机制}

口腔医学学习内容包含疾病定义、诊断条件、治疗原则、用药说明和操作要求，术语、数值或适用条件发生偏移都会改变医学含义。OralSage 将检查环节安排在输入接收、资料检索、上下文构造、内容生成和结果输出五个阶段，使用户请求、检索材料、生成陈述与引用位置保持明确对应。

\subsection{多阶段医学内容处理}

资料进入知识库时已进行 PDF 版面和网页层级恢复，清理 OCR 异常、重复页眉、模板文本及参考文献密集片段，并保留文档、页面、片段和标题路径。重复文本、免责声明和低信息内容写入质量标记，知识节点与文本之间的链接经过标识、目标和位置检查后参与检索。由此，模型接收的医学内容已经完成基础清理，并保留能够定位到具体资料的来源信息。不同资料按照内容类型进入生成过程。系统优先采用临床规范与教材中的明确结论，同时保留资料类型，使通用医学参考和考试材料在各自范围内参与内容组织。

在线检索阶段由 Quality Gate 读取来源属性与质量标记，过滤模板化内容，降低重复文本和参考文献密集片段的排序优先级。Alias、BM25、BGE-M3 和 Graph-PPR 返回的候选按照统一对象标识合并，再由 RRF 形成综合顺序。多路结果共同进入排序，可以减少单一关键词相似或单一语义相似带来的偏差。

\subsection{检索约束和引用绑定}

KAGSolver 将医学文本、知识节点、关系路径、资料位置和检索记录进行上下文统一，并要求医学陈述使用上下文中的材料，限制自行补充外部知识，并对药物剂量、诊断阈值、禁忌证、操作条件和疾病分类采用关键术语保留策略。

引用信息按照以下位置链组织，使生成内容能够定位到具体资料：

\begin{center}
  \softid{智能体上下文 -> 检索文本单元 -> 来源片段 -> 页面 -> 资料}
\end{center}

Citation Builder 为返回内容补充资料标题、页码范围和局部片段，KG Path Builder 记录查询种子与相关知识对象之间的关系路径，Retrieval Trace 保存召回通道、候选数量和融合顺序。通过引用绑定，使回答中的医学陈述与相应文本位置一并输出，并为后续审查提供可核对的内容单元\cite{mirage2024}。

Resource Factory 在生成讲义前先从检索结果中选择本节所需的知识节点、知识关系、题目和文本单元，再执行定向检索补充材料。内容生成完成后，程序核对实际使用的对象标识是否属于本次检索允许集合，定向检索获得的文本按照文本单元参与引用，知识节点和关系则保持知识库中已有的对象类型，从而避免生成过程自行构造不存在的知识对象或资料位置。

\subsection{输出过滤与分级审查}

输出审查首先检查结构完整性。后端依据预定义数据结构核对字段类型、必填内容、文本长度和对象标识，无法解析的结构不会进入资源派生与页面展示。对于带有医学陈述的讲义和回答，程序继续检查引用标识是否存在、引用位置是否属于相应资料，以及生成结果使用的知识节点、关系和题目是否来自本次允许集合。

系统重点核对疾病名称、分型、药物剂量、诊断阈值、适应证、禁忌证和操作条件。一般讲解内容检查陈述与引用文本的语义对应，包含数值和临床条件的内容还需比对关键术语与数值。对于需要细粒度复核的回答，审查程序将回答拆分为独立医学陈述，并使用 BGE Reranker 计算陈述与检索材料的对应分数，随后标记需要重新检索或改写的内容，保证医学一致性。

内容安全检查面向个体化诊断、处方剂量、侵入性操作、自伤伤人、药物滥用和其他可能造成现实伤害的输出。检测到高风险内容后，系统移除可直接执行的步骤、剂量和操作细节，将回答转换为医学常识、风险提示、就医建议或紧急求助指引。面向课程学习的定义、机理和规范性知识继续保留，输出范围限定在教学解释和经过引用的医学材料内。

\subsection{有害输入识别}

当输入包含要求忽略既有规则、泄露系统提示词、伪造引用、虚构资料位置或改变智能体职责的指令时，处理程序保留其中可识别的学习主题，阻断其控制意图。检索获得的网页文字和题库文本同样按医学材料处理，其中出现的命令式语句不会替代系统指令。

输入审查将请求区分为正常学习、医学高风险、指令操纵和有害行为等类型。涉及自伤伤人、毒物使用、处方药滥用、危险器械操作或规避医疗监管的请求进入安全响应流程，系统不提供可执行步骤、精确剂量和规避方法，转而给出风险说明、专业医疗建议和紧急求助信息。经过拆字、编码、同义替换或多轮诱导重复提交的请求仍按规范化后的语义进行判断，保持输入处理规则的一致性。
